In the general case, $G$ will be isomorphic to a closed subscheme of a scheme of the form $S[(t_i)]$, \ie $C$ will be a quotient of an $A$-algebra of the form $A[(t_i)]$. Let $(F_j)$ be a system of generators of the ideal by which one takes the quotient. Suppose $\mathscr B$ free of rank $n$, which is permissible on covering $S$ by smaller affine open subsets. Choose a basis $(e_k)_{1\le k\le n}$ of $\mathscr B$; then, writing the $n$ components with respect to this basis of $F_j((x_i))$, for $x_i=\sum x_{ik}e_k$ ($x_{ik}$ indeterminate coefficients, taken in an unspecified algebra $A'$ over $A$), one finds, for every $F_j$, $n$ polynomials $F_{j,k}$ in the $(x_{i,k})_{i\in I,\,1\le k\le n}$, with coefficients in $A$. One observes at once that $\uHom_S(H,G)$ is represented by the spectrum of the quotient of the polynomial ring $A[(x_{i,k})]$ by the ideal generated by the $F_{j,k}$. This proves b) at once.

Let us prove a). For every ordinary finite commutative group $M$, let $F_M$ be the subfunctor of $F$ obtained by restricting oneself to subgroups of $G_T$ that are of multiplicative type and of type $M$. One sees easily, by a gluing argument like the one used in 3.5 (which we should have stated by breaking things down a little more!), that it suffices to verify that the $F_M$ are representable; then $F$ will be representable by the sum prescheme of the $F_M$, where $M$ runs through the set of classes of finite commutative groups up to isomorphism. ($F$ is indeed the sum of the $F_M$ in the category of sheaves$\ldots$).

Henceforth we suppose $M$ fixed, and shall write $F$ instead of $F_M$. Let $H=D_S(M)$, consider the subfunctor $F'=\underline{\mathrm{Imm}}_{S\text{-gr}}(H,G)$ of $\uHom_{S\text{-gr}}(H,G)$ whose value at $T$ is the set of homomorphisms of $T$-groups $H_T\to G_T$ that are closed immersions. One already knows by b) that $\uHom_{S\text{-gr}}(H,G)$ is representable by an affine $S$-prescheme, and using IX~6.8 one sees at once that $F'$ is representable by an open and closed sub-prescheme of the latter, hence is also affine over $S$. Finally, consider the canonical morphism $F'\to F$, which associates to each monomorphism $H_T\to G_T$ the image group. One observes at once, by the definitions, that in this way $F'$ becomes a principal homogeneous bundle (in the category of sheaves for fpqc) over $F$, with group $\Gamma_F=\Gamma_S\times_S F$, where $\Gamma=\Aut_{\mathrm{gr}}(M)$, whence it follows ``by descent'' that this morphism is representable (\ie for every morphism $T\to F$, with $T$ representable, $T\times_F F'$ over $T$ is representable---in fact, representable by a principal Galois bundle under $\Gamma$). Thus, by IV~4.6.6, $F$ is representable if and only if the quotient $F'/F''$, where $F''=F'\times_F F'$, exists in $\Sch$ and is universally effective for faithfully flat quasi-compact morphisms, or, what amounts to the same thing, if and only if the quotient $F'/\Gamma$ exists and is universally effective for the said morphisms. Now, since $F'$ is affine over $S$, we saw in V~4.1 that the condition in question is indeed satisfied. This proves the representability of $F$ in a).

As for the addition, relating to the case where one supposes $G$ of finite presentation over $S$, it follows at once from the preceding proof, taking into account that by b), $F'$ is then locally of finite presentation over $S$, so that one can apply Exposé~V\nde{Reference to be checked/made precise.}.

\sgasection{4}{The scheme of subgroups of multiplicative type of a smooth affine group}
The main result of the present exposé consists of the
\begin{theorem}{4.1}\label{XI.4.1}
Let $S$ be a prescheme, $G$ an $S$-group prescheme smooth and affine over $S$, $F$ the functor $\Sch^\circ_{/S}\to\Ens$ defined by
\[
F(T)=\text{set of subgroups of multiplicative type of }G_T.
\]
Then the functor $F$ is representable, and is smooth and separated over $S$.
\end{theorem}
Let us point out at once the following variant:
\begin{corollary}{4.2}\label{XI.4.2}
Let $G,H$ be two $S$-group preschemes, with $G$ smooth and affine over $S$, $H$ of multiplicative type and of finite type. Then $\uHom_{S\text{-gr}}(H,G)$ is representable, and is smooth and separated over $S$.
\end{corollary}
Indeed, when $H$ is smooth over $S$, so is $H\times_S G$, and one can apply 4.1 to the latter. By considering graph groups, $\uHom_{S\text{-gr}}(H,G)$ becomes a subfunctor $F'$ of the functor $F$: ``subgroups of multiplicative type of $H\times_S G$'', namely $F'(T)=$ set of subgroups $K$ of multiplicative type of $(H\times_S G)_T=H_T\times_T G_T$ such that the homomorphism induced by the first projection
\[
K\to H_T
\]
is an isomorphism. By IX~2.9 the canonical morphism $F'\to F$ is an open immersion, so, since $F$ is representable, so is $F'$, which will be representable by an open subset of $F$; and since $F$ is separated and smooth over $S$, so will $F'$ be. In the case where $H$ is finite over $S$, it suffices to apply 3.12 b) for the representability of $\uHom_{S\text{-gr}}(H,G)$. When $H$ is a product $H_1\times_S H_2$, with $H_1$ smooth over $S$ and $H_2$ finite over $S$, then $\Hom_{T\text{-gr.}}(H_T,G_T)$ is identified with the subset of $\Hom_{T\text{-gr.}}((H_1)_T,G_T)\times\Hom_{T\text{-gr.}}((H_2)_T,G_T)$ formed by pairs $(u_1,u_2)$ such that $u_1\cdot u_2=u_2\cdot u_1$, whence it follows that $\uHom_{S\text{-gr}}(H,G)$ is representable by a closed sub-prescheme of $\uHom_{S\text{-gr}}(H_1,G)\times_S\uHom_{S\text{-gr}}(H_2,G)=X$, as one sees by applying VIII~6.5 b), where one takes $Y=H$, $Z=G$, $q_1$ and $q_2$ being defined respectively by $(u_1,u_2)\mto u_1\cdot u_2$ and $(u_1,u_2)\mto u_2\cdot u_1$.

In the general case, since the question is local on $S$ for the Zariski topology, one can suppose $S$ affine, and $H$ of constant type over $S$. Then, since $H$ is quasi-isotrivial (X~4.5), one can find an étale surjective morphism $S'\to S$, $S'$ affine, which splits $H$, \ie such that $H'=H_{S'}$ is diagonalizable. Then the preceding result applies, for a diagonalizable group is the product of a diagonalizable torus with a finite diagonalizable group. It remains to see that the descent datum obtained on the $S'$-prescheme $\uHom_{S'\text{-gr}}(H',G')=X'$ is effective. This is seen by the argument of X~5.4, noting that in \loccit, the hypothesis that $X'\to S'$ was separated and locally quasi-finite served only to ensure that every open subset of $X'$ quasi-compact over $S'$ was quasi-affine over $S'$. Now this property is still satisfied in the present case, as follows easily from the fact that it holds for the functor $F$ of 4.1 (which will be seen in the course of the proof of 4.1). This (or any other variant of this little reduction) establishes the representability of $\uHom_{S\text{-gr}}(H,G)$, and at the same time the fact that it is separated over $S$. It is locally of finite presentation over $S$, as one sees, for example (as was pointed out in 3.2), thanks to the fact that this functor ``commutes with inductive limits of rings''. Finally, since this functor is formally smooth over $S$ (by 2.1), it is smooth over $S$.

\begin{remark}{4.3}\label{XI.4.3}
We have here deduced 4.2 from 4.1, which is really immediate only when $H$ is also smooth over $S$. For the deduction to be made without contortions in the general case, the representability result 4.1 would have to be established without supposing $G$ smooth over $S$, but only affine of finite presentation over $S$. (Of course, then $F$ will in general no longer be smooth over $S$!). There is little doubt that 4.1 remains true under these more general hypotheses, but the proof seems to have to be more delicate (for lack of being able to invoke 3.8){\renewcommand{\thefootnote}{*}\footnote{It is indeed proved for $G$ flat and quasi-affine over $S$ with connected fibers, provided one restricts oneself to central subtori of $G$ (XV~8.8).}\addtocounter{footnote}{-1}}\nde{In the case where $H$ is smooth (not necessarily affine) over a locally noetherian normal base $S$, M.~Raynaud showed that the largest representable open subset of the functor of subtori of $H$ is a disjoint sum of smooth affine open subsets over $S$. This is theorem IX.9.26 in \textit{Faisceaux amples sur les schémas en groupes et sur les espaces homogènes}, Lecture Notes Maths.~119 (1970).}. Let us point out, however, that when $G$ is a closed subgroup of an affine smooth group $G'$ over $S$, then the functor $F$ representing subgroups of multiplicative type of $G$ is representable by a closed sub-prescheme of the prescheme representing the analogous functor $F'$ for $G'$ (falling under 4.1), as one sees easily by applying VIII~6.4. This also raises the question: is a group scheme $G$ over an affine $S$, which is affine and of finite presentation over $S$, isomorphic to a group subscheme of a $\GL(n)_S$, for suitable $n$? This is true when $S$ is the spectrum of a field, \cf VI$_{\mathrm B}$~11.11, but unfortunately false in general, even for tori, \cf 4.6. Finally, note that one could also prove 4.2 directly by exactly the same method as 4.1.
\end{remark}

Let us now prove 4.1. Since the functor $F$ is obviously of a local nature, one can suppose $S$ affine, hence $S=\Spec(A)$, $A$ a ring. Considering $A$ as the inductive limit of its subrings of finite type over $\ZZ$, and noting that $G$ comes from a smooth affine group over such a subring (EGA~IV~8), one is reduced to the case where $S$ is noetherian. For every integer $n>0$, let $T_n$ be the functor defined like $F$, but restricting oneself to subgroups $H$ of multiplicative type of $G_T$ such that $n\cdot\mathrm{id}_H=0$, \ie such that ${}_nH=H$. Order the set $I$ of integers $>0$ by the divisibility relation. When $m$ is a multiple of $n$, define
\[
p_{n,m}:T_m\to T_n
\]
by $p_{n,m}(H)={}_nH=\Ker(n\cdot\mathrm{id}_H)$.

In this way the $T_n$ form a projective system of functors over $S$. By 3.12 a), the functors $T_n$ are representable, affine and of finite type over $S$. Define likewise morphisms
\[
u_n:F\to T_n,
\]
by the relation $u_n(H)={}_nH$. In this way one obtains a morphism
\[
u:F\to T=\varprojlim_n T_n,
\]
where the $\varprojlim$ is taken in the category of functors over $S$. But let us point out at once that, since the $T_n$ are representable and affine over $S$, so is $T$ (it will be the spectrum of the inductive limit of the quasi-coherent algebras over $S$ defining the $T_n$). Of course $T$ is not in general of finite type over $S$.

We shall apply 3.7 and are reduced to verifying conditions a) to d) of 3.7, which will imply that $F$ is representable by a prescheme locally of finite type over $S$. It then follows from 2.1 bis that $F$ is even smooth over $S$, and since $F$ is a subfunctor of $T$, which is affine over $S$, it follows that $F$ is separated over $S$ (being separated over $T$, which is separated over $S$). Let us prove at once the addition invoked above, namely that every open subset $U$ of $F$ quasi-compact over $S$ is quasi-affine over $S$: this follows from 3.11 and the fact that the $T_u$ are affine over $S$.
% typo?: T_u kept as printed.

Let us therefore verify the conditions of 3.7.

a) $F$ is a sheaf for the faithfully flat quasi-compact topology, by descent theory SGA~1 VIII, which applies here since groups of multiplicative type over $T$ are affine over $S$. It commutes with inductive limits of rings by the general groundwork EGA~IV~8. Let us show that it commutes with adic projective limits of artinian local rings. When one is dealing, instead of the functor $F$, with the analogous functor considered in 4.2, this property is nothing other than that of IX~7.1 in the special case where $A$ is a complete noetherian local ring, endowed with an ideal of definition for its usual topology (N.B. It is precisely here that the hypothesis $G$ affine enters essentially). In the present case, we are reduced to proving the
\begin{lemma}{4.4}\label{XI.4.4}
Let $A$ be a complete noetherian local ring, with maximal ideal $\mathfrak m$, $G$ an affine group scheme over $S=\Spec(A)$. For every integer $n\ge0$, let $S_n=\Spec(A/\mathfrak m^{n+1})$, $G_n=G\times_S S_n$. For every $n$, let $H_n$ be a subgroup of multiplicative type and of finite type of $G_n$, such that for $m\ge n$, $H_n$ is deduced from $H_m$ by reduction. Under these conditions, there exists a unique subgroup $H$ of multiplicative type in $G$ that reduces to the $H_n$.
\end{lemma}
By X~3.2 there exists a group $H$ of multiplicative type over $S$, necessarily of finite type and isotrivial, determined up to unique isomorphism, endowed with an isomorphism $H\times_S S_0\simeq H_0$ (using the fact that $H_0$ is isotrivial, being of finite type over a field, \cf X~1.4). By X~2.1, for every $n$, the isomorphism $H\times_S S_0\simeq H_0$ lifts to a unique isomorphism $H\times_S S_n\simeq H_n$. This being said, by IX~7.1 already cited, the homomorphisms $H_n\to G_n$ come from a unique homomorphism of $S$-groups $u:H\to G$.

By IX~6.6 the latter is a monomorphism since $u_0:H_0\to G_0$ is. This completes the proof of 4.4.

b) The morphism $u:F\to T$ is a monomorphism. This follows from the density theorem IX~4.7, in the form of corollary 4.8 b). One should take care that it is essential, for the application we make of it here, to have this result over a base that is not necessarily noetherian. (N.B. Since the functor $T$ does not commute with inductive limits of rings, it is not possible to reduce to that case a priori).

b') The morphism $u:F\to T$ is infinitesimally étale, in other words one has the
\begin{lemma}{4.5}\label{XI.4.5}
Let $A$ be an artinian local ring with residue field $k$, $S=\Spec(A)$, $I$ an ideal $\subset\operatorname{rad}(A)$, $S'=\Spec(A/I)$, $G$ an $S$-group prescheme smooth over $S$, $G'=G\times_S S'$; give ourselves a subgroup $H'$ of multiplicative type of $G'$ and, for every integer $n>0$, a subgroup $H(n)$ of multiplicative type of $G$, in such a way that
\begin{enumerate}[label=\arabic*$^\circ$)]
\item for every multiple $m$ of $n$, $H(n)={}_nH(m)$, and
\item $H(n)'={}_nH'$.
\end{enumerate}
Under these conditions, there exists one and only one subgroup $H$ of multiplicative type of $G$ such that ${}_nH=H(n)$ for every $n$.
\end{lemma}
Uniqueness is already contained in b). For existence, an immediate induction reduces us to the case where $J\mathfrak m=0$, $\mathfrak m$ being the maximal ideal of $A$. Let $k=A/\mathfrak m$, $S_0=\Spec(k)$, $G_0=G\times_S S_0$, $\mathfrak g_0$ the Lie algebra of $G_0$, $\mathfrak h_0$ that of $H_0$. One has a canonical group isomorphism:
\[
\mathfrak g_0\otimes_k J\simeq\Ker(G(S)\to G(S_0)),
\]
\cf Exp.~III. By IX~3.6 bis and 3.7, there exists a subgroup $H$ of multiplicative type of $G$ reducing to $H'$, and such an $H$ is determined modulo inner automorphism by an element of $N=\Ker(G(S)\to G(S_0))$. Thus the set $P$ of liftings $H$ of $H_0$ is a principal homogeneous set under the group $N/M$, where $M$ is the group of $g\in N$ such that $\operatorname{int}(g)\cdot H=H$.
% typo?: J used after the statement defines I; kept as printed.

Now one sees easily (Exp.~III) that this subgroup is nothing other than the vector subspace of $\mathfrak g_0\otimes_k J$ formed by the invariants under $H_0$, when $H_0$ acts on $\mathfrak g_0\otimes_k J$ by the representation induced by the adjoint representation of $G_0$. Likewise, the set $P(n)$ of liftings of $H(n)_0={}_nH_0$ is a principal homogeneous set under $N/M(n)$, where $M(n)$ is the vector subspace of $\mathfrak g_0\otimes_k J=N$ formed by the invariants under ${}_nH_0$. Using the density theorem Exp.~IX~4.7, one sees easily that for large $n$ (in the sense of the order relation put on the set of integers $n>0$, namely the divisibility relation) one has $M=M(n)$. Consequently the natural map $H\mto{}_nH$ from $P$ to $P(n)$, which is compatible with the actions of $N$ and hence with the homomorphism $N/M\to N/M(n)$ on the groups of operators, is bijective for large $n$. The conclusion of 4.5 follows at once.

To verify c) and d) of 3.7, we use 3.8, which reduces us to the verification of c') and d') below.

c') The $T_n$ are smooth over $S$, and the transition morphisms $p_{n,m}:T_m\to T_n$ are smooth.

This is nothing other than 2.2 bis.

d') With the notation of 3.8, a point $\xi$ of $F$ with values in a field $k$ over $S$ is nothing other than a subgroup $H_0$ of multiplicative type of $G_0=G_k$. Taking up again the argument of b') above, one sees that the integer $d(\xi)$ considered in 3.8 is nothing other than the dimension of $\mathfrak g_0/\mathfrak g_0^{H_0}$, where $\mathfrak g_0$ is the Lie algebra of $G_0$ and $\mathfrak g_0^{H_0}$ is the vector subspace of invariants under $H_0$. It is therefore a finite integer, \ie condition d') $1^\circ$) of 3.8 holds. With the notation of d') $2^\circ$, giving a morphism $v:X\to F$ amounts to giving a subgroup $H$ of multiplicative type of $G_X$. For $x\in X$, the integer $d(\xi_x)$ is then nothing other than the dimension of $(\mathfrak g\otimes\kappa(x))^{H_x}$, where $\mathfrak g$ is the sheaf of Lie algebras of $G_X$ (which is locally free of finite type on $X$, for $G_X$ is smooth over $X$, and $\mathfrak g\otimes\kappa(x)$ is nothing other than the Lie algebra of the fiber $G_x$ of $G_X$ at $x$), and $H_x$ is the fiber of $H$ at $x$. Now, since $\mathfrak g$ is a locally free module over $S$ on which the group $H$ of multiplicative type acts, one sees at once that the subfunctor $\mathfrak g^H$ of invariants under $H$ is given by a subsheaf locally a direct summand, hence locally free, of $\mathfrak g$. (By descent, one is reduced to the case where $H$ is diagonalizable, and where one applies Exp.~I~4.7.3, noting that the subsheaf of invariants corresponds to the component of degree zero). Consequently $d(\xi_x)=\text{rank at }x\text{ of }\mathfrak g^H$, hence it is a locally constant function of $x$. This completes the proof of condition d').
% typo?: the quotient dimension becomes the invariant dimension in this paragraph; retained as printed, including the base S.

We have thus verified the conditions of 3.7, which completes the proof of 4.1.

\begin{remark}{4.6}\label{XI.4.6}
When $G=\GL(n)_S$, one can give a direct proof of 4.1 that is distinctly simpler and more explicit, using I~4.7.3. The proof shows in addition that, in this case, the moduli scheme is a sum scheme of a family of affine schemes over $S$. Proceeding as was said in 4.3, one deduces the same result whenever $G$ is a closed subgroup of a group of the form $\GL(n)_S$. One should beware, however, of believing that the preschemes representing the functors in 4.1 or 4.2 are always sums of a family of affine schemes over $S$. Let, for example, $H$ be a group of multiplicative type and of finite type over a locally noetherian prescheme $S$; then by X~5.11, $H$ is isotrivial if and only if $\uHom_{S\text{-gr}}(H,\Gm)$ is a sum of a family of affine preschemes over $S$, but it was pointed out (X~1.6) that non-isotrivial tori $H$ (of relative dimension 2) can exist; for such a torus, $\uHom_{S\text{-gr}}(H,\Gm)$ is therefore not a sum of affine $S$-preschemes over $S$, and one sees likewise that the ``dual'' twisted constant group $\uHom_{S\text{-gr}}(\Gm,H)$ is not such a sum either (for if two commutative twisted constant groups of finite presentation $R,R'$ are dual to each other, $R'=\uHom_{S\text{-gr}}(R,\ZZ_S)$, one sees easily that one is isotrivial if and only if the other is). This last point also shows that such an $H$ is not isomorphic to a subgroup of a group of the form $\GL(n)_S$; more precisely, one can show that a group of multiplicative type and of finite type over a connected locally noetherian $S$ is isomorphic to a subgroup of a group of the form $\Aut_{\Modcat}(\mathscr E)$ (with $\mathscr E$ a locally free module of finite type over $S$) if and only if it is isotrivial. Finally, taking $G=H\times\Gm$ in the two preceding examples, one finds an example where the prescheme representing the functor $F$ of 4.1 is not a sum of a family of affine $S$-preschemes over $S$ (with $G$ a torus of relative dimension 3, if desired).
\end{remark}
